% !TeX root = ../main.tex

\chapter{绪论}
\label{ch:intro}

\section{研究背景}

遥感图像中的船舶检测是海上监管、港口调度、海事执法与应急救援中的基础能力之一。与自然场景目标检测相比，遥感图像通常覆盖范围更大、背景更复杂、目标更稀疏，且尺度变化明显。纯人工判读在效率、稳定性和可追溯性上都存在明显限制。因此，把目标检测技术应用到遥感船舶场景，并形成可复现、可演示、可回放的系统链路，具有明确的工程价值。

近几年，通用目标检测方法从 two-stage、one-stage 持续发展到 anchor-free 与 transformer 检测器，精度和推理效率都在提升\cite{ren2015_fasterrcnn,lin2017_fpn,lin2017_focalloss,tian2019_fcos,redmon2016_yolo,carion2020_detr,zhu2021_deformabledetr}。与此同时，遥感小目标检测研究也在持续强调尺度变化、背景干扰和样本稀疏等问题对模型稳定性的影响\cite{wang2023_sod_review,zhao2024_shipsurvey}。但在中分辨率遥感图像中，微小船舶目标像素占比很小，又容易受到海浪、云层、岸线和近岸建筑干扰，通用检测器直接迁移时常出现漏检和误检。

随着公开数据和开源工具链逐步完善，研究重点也从“能不能训起来”转向“能不能统一评测、稳定复现并系统化落地”。仅完成单模型训练难以支撑完整结论，还需要在统一协议下比较不同检测范式，并把算法能力落实到可运行系统中。

\section{研究问题与挑战}

本课题聚焦“遥感场景下的微小船舶检测系统设计与实现”。核心问题是：在复杂背景和弱目标特征条件下，如何平衡检测精度、召回能力与工程可用性。结合 LEVIR-Ship 数据集特点及已有研究\cite{chen2022_tinyship_drenet}，主要挑战可归纳为四点。

第一，\textbf{目标尺度小、纹理弱}。微小船舶在图像中只占很少像素，下采样后特征容易被淹没。第二，\textbf{背景干扰强}。海浪、云层、岸线和港口设施会产生大量类目标纹理。第三，\textbf{样本分布不均}。不同场景目标密度差异较大，正负样本比例不平衡。第四，\textbf{评测与复现口径不统一}。数据划分、输入策略、阈值和评测脚本差异会直接影响横向比较可信度。

因此，本课题不只关注单一指标提升，而是同时推进统一实验协议、多模型比较和系统实现，保证结论可解释且可复核。

\section{研究目标与研究内容}

课题总体目标是：以 LEVIR-Ship 为统一实验基准，完成多模型对比实验、结果分析与检测系统实现，形成“问题分析、实验论证、系统落地、论文表达”闭环。

围绕该目标，主要工作包括：

\begin{enumerate}
  \item \textbf{统一数据与评测协议}：统一数据划分、评价指标和结果记录格式，建立可复验流程；
  \item \textbf{完成多模型对比实验}：选择 DRENet、FCOS 和 YOLO26，开展训练、测试与横向比较；
  \item \textbf{开展误差与消融分析}：围绕复杂背景、弱目标和阈值策略分析成功检测、漏检与误检；
  \item \textbf{实现轻量级检测系统}：基于 React、FastAPI、SQLite 完成多模型接入、异步任务执行与结果回看。
\end{enumerate}

\section{技术路线与主要工作}

技术路线如图~\ref{fig:ch1_pipeline} 所示，整体遵循“数据准备 $\rightarrow$ 模型训练与推理 $\rightarrow$ 结果分析 $\rightarrow$ 系统集成”。与只做算法训练的路径不同，这里把实验链路和系统链路并行推进，让离线结果和在线演示保持同口径。

\begin{figure}[htbp]
  \centering
  \includegraphics[width=0.96\textwidth]{generated/ch1_pipeline.png}
  \caption{技术路线示意}
  \label{fig:ch1_pipeline}
\end{figure}

从图~\ref{fig:ch1_pipeline} 可见，上半部分是实验论证链路，下半部分是系统证据链路。两条链路在统一数据协议和统一推理接口处汇合，保证论文中的表格、图示与系统任务回放一致。

基于这一路线，已形成三项核心产出：一是统一实验协议；二是 DRENet、FCOS、YOLO26 的对比结果与误差分析；三是复用现有推理内核实现 Web 检测系统，打通任务提交、模型推理、结果存储和前端展示。

\section{论文结构安排}

后文按“相关工作、方法设置、实验分析、系统实现、总结展望”展开：

\begin{enumerate}
  \item 相关工作部分梳理通用检测、小目标检测与遥感船舶检测研究，并明确本文切入点；
  \item 方法设置部分介绍数据集、评价指标、模型方案和统一实验协议；
  \item 实验分析部分给出三模型对比、效率结果、定性样例与消融分析；
  \item 系统实现部分说明需求、架构、数据库、接口和关键流程；
  \item 总结展望部分归纳主要结论，并说明当前边界与后续优化方向。
\end{enumerate}

\section{小结}

本部分明确了课题背景、核心挑战、研究目标和整体路线。后文将围绕统一评测协议展开相关研究综述、实验结果分析和系统实现说明，形成完整论证链路。
